% Capítulo 1: Introducción
% Propósito: Describir los antecedentes, planteamiento del problema, justificación, formulación y los objetivos (general y específicos) del proyecto de grado.


\section{Introducción}

\section{Planteamiento del proyecto}
Actualmente, la gestión que tienen los reportes de incidencias comunitarias en la Alcaldía Municipal de Rivas presenta limitaciones relacionadas con la ausencia de un canal digital centralizado que permita a los ciudadanos del municipio integrar su recepción, organización y seguimiento. La utilización de diferentes medios para comunicar las problemáticas de la comunidad puede ocasionar que la información se encuentre fragmentada y que no exista una estructura uniforme para registrar los datos asociados a cada incidencia. Asimismo, se presentan limitaciones para incorporar desde el momento del reporte información complementaria, como la ubicación geográfica y evidencia fotográfica, elementos que pueden facilitar la identificación y posterior atención de las incidencias comunicadas por los ciudadanos.

Estas condiciones dificultan la gestión de las incidencias durante las distintas etapas de su atención. Al no disponer de un mecanismo centralizado que concentre la información de cada reporte, pueden presentarse dificultades para mantener un seguimiento continuo y uniforme de los casos, así como para conocer con claridad su estado durante el proceso de atención. Esto puede contribuir a prolongar los tiempos necesarios para canalizar las incidencias hacia las áreas correspondientes y dificulta mantener un registro organizado de las acciones realizadas sobre cada caso.

Esta situación afecta tanto a los ciudadanos que realizan reportes de incidencias comunitarias como al personal municipal encargado de gestionarlos. Desde la perspectiva ciudadana, las limitaciones en el seguimiento dificultan conocer si una incidencia ha sido recibida, se encuentra en proceso de atención o ha sido resuelta, reduciendo la retroalimentación disponible durante el proceso. Por otra parte, para el personal de la alcaldía, la información dispersa y poco estructurada dificulta las tareas de consulta, clasificación y seguimiento de los reportes, especialmente cuando se requiere revisar información correspondiente a incidencias registradas con anterioridad.

La ausencia de información centralizada también limita las posibilidades de aprovechar los registros históricos de incidencias como apoyo para la gestión municipal. Disponer de información organizada permitiría identificar los tipos de problemas que se presentan con mayor frecuencia, reconocer sectores donde existe una concentración recurrente de reportes y generar métricas que faciliten el análisis de las necesidades de la comunidad. Este aspecto resulta relevante para la planificación y priorización de las intervenciones municipales, debido a que la información recopilada puede transformarse en un recurso para apoyar la toma de decisiones y orientar la atención hacia las zonas que presentan una mayor constancia de incidencias.

Ante esta problemática, se plantea el diseño y desarrollo de una plataforma web orientada a centralizar el envío, clasificación y seguimiento de incidencias comunitarias en la Alcaldía Municipal de Rivas. La plataforma permitiría que los ciudadanos registren sus incidencias mediante información estructurada que incluya una descripción sobre la incidencia, evidencia fotográfica y ubicación geográfica, mientras que el personal municipal dispondría de herramientas para gestionar y actualizar el estado de cada caso. Asimismo, la incorporación de mecanismos de visualización geográfica permitiría representar las incidencias sobre un mapa e identificar sectores con mayor concentración de reportes, complementando la gestión operativa con información que pueda ser utilizada para el análisis y la toma de decisiones municipales.

\section{Objetivos}

\textbf{Objetivo General}

Diseñar una plataforma web para la gestión de reportes de incidencias comunitarias 
para la alcaldía del municipio de Rivas, Nicaragua, con el propósito de optimizar los 
procesos de respuesta y recepción de peticiones ciudadanas en digital.

\textbf{Objetivos Específicos}
\begin{itemize}
	\item Analizar los requerimientos funcionales y no funcionales del sistema, 
	diagnosticando el proceso actual y las principales deficiencias en los canales de 
	comunicación entre los ciudadanos y la alcaldía municipal.
	
	\item Diseñar la arquitectura del software, el modelo de base de datos y las interfaces 
	de usuario, estructurando una plataforma que permita el registro de incidencias 
	con evidencias fotográfica y ubicación geográfica de manera digital.
	
	\item Desarrollar los módulos funcionales de la aplicación, incluyendo un panel de 
	administración que provea métricas y actualización de estados para elevar la 
	transparencia y eficiencia en la atención de incidencias.
	
	\item Realizar pruebas funcionales, de usabilidad y de seguridad en el sistema para 
	garantizar su correcto desempeño y validar el cumplimiento de los 
	requerimientos establecidos.
	
	\item Elaborar la documentación técnica y los manuales de usuario de la plataforma, 
	creando una guía clara que sirva de base para la capacitación del personal y la 
	comunidad.
\end{itemize}

\section{Antecedentes}

A nivel global, la adopción de herramientas tecnológicas ha mejorado la gestión pública y la resolución de problemáticas locales, agilizando procesos con precisión. En este contexto, la sistematización del registro y seguimiento de incidencias comunitarias en [Nombre de tu ciudad/municipio] representa una oportunidad para mejorar la atención ciudadana y la comunicación digital entre ciudadanos e institución. A continuación, se presentan investigaciones previas que fundamentan el desarrollo de este trabajo.
Ámbito Internacional:
En el ámbito internacional(\cite{parrales2025desarrollo}) desarrolló la investigación titulada “Desarrollo de un sistema web para el reporte de daños en la infraestructura urbana de la parroquia Anconcito integrando Google Cloud Vision API”, presentada en la Universidad Estatal Península de Santa Elena (Ecuador). 
El estudio abordó las deficiencias en la gestión y atención de fallas urbanas en el barrio Luis Celleri.
Para mitigar esta problemática, el autor diseñó e implementó una plataforma web responsiva bajo una arquitectura cliente-servidor, empleando Angular en el frontend, Laravel en el backend y MySQL como base de datos.

El sistema permite a los ciudadanos registrar reportes detallados incorporando coordenadas de geolocalización precisa e imágenes descriptivas. 
Según esta investigación la Implementación del sistema digital con geolocalización y mapas interactivos para el reporte ciudadano de infraestructura urbana dio beneficios sólidos sustentados con evidencia los cuales son: 

Reducción del tiempo en la creación de reporte: El 86.7\% de los usuarios completó sus reportes en menos de diez minutos, lo que representa una reducción del 89\% en el tiempo necesario en comparación con los métodos tradicionales. 
Eliminación de barreras del sistema tradicional: La plataforma eliminó la necesidad de realizar desplazamientos físicos (identificado como una traba prioritaria por el 60\% de los participantes) y redujo el papeleo excesivo (señalado por el 53.3\%).
 
Este trabajo aporta a la presente investigación un precedente del uso de coordenadas geográficas en plataformas de monitoreo ciudadano, sirviendo de referencia en nuestro proyecto. Además, el conocimiento de funcionalidades tales como: herramientas de geolocalización, mapas interactivos y que la implementación de estados de seguimiento en los reportes es crucial en la bidireccionalidad confiable entre las autoridades locales y la ciudadanía fomentando un modelo de gestión colaborativa.

En el Municipio de Othón P. Blanco, en el estado de Quintana Roo, México se diseñó e implemento una plataforma que sustituyo al antiguo sistema de atención formado únicamente por llamadas telefónicas. " 072 Reportes”(\cite{trejo2019sistema}) surgió

con el objetivo principal de optimizar, agilizar y transparentar la recepción, canalización, monitoreo y solución de las fallas en los servicios públicos municipales (como baches, alumbrado público deficiente, coladeras tapadas o mantenimiento de parques). Está conformada por:

Aplicación móvil (Android) para la ciudadanía: Permite a los habitantes levantar reportes capturando la ubicación geográfica exacta vía GPS, agregando una descripción y adjuntando hasta tres fotografías del desperfecto. También permite a los usuarios recibir notificaciones del avance de su reporte, consultar avisos oficiales y responder encuestas de satisfacción de servicio.
Y un Sistema web de administración (Intranet) para el ayuntamiento: Utilizado exclusivamente por el personal municipal (administradores, gestores y áreas operativas responsables) para clasificar las incidencias, asignar tareas de atención, registrar avances, visualizar mapas de localización de reportes y consultar gráficas estadísticas en tiempo real.

La implementación de este sistema app y web aporto al municipio la modernización de sus procedimientos de recepción, seguimiento y solución de reportes de servicios públicos mediante la integración de las TIC, y la computación móvil por medio de esta plataforma digital, el ayuntamiento logra optimizar sus procesos de atención y coordinación entre áreas, resolver los errores de precisión en la ubicación geográfica de los problemas y evitar el desperdicio de recursos humanos, tiempo y combustible por motivo de reportes falsos, proyectando de este modo una imagen institucional moderna y confiable ante la población.

Esta investigación aporta un catálogo estructurado de requerimientos funcionales y el diseño de diagramas entidad-relación para bases de datos los cuales serán útiles al reconocer la estructura que llevan este tipo de sistemas.

\section{Justificación}
% Revisar la de la carta

\section{Alcance}
El alcance de la presente investigación se define metodológicamente como descriptivo y propositivo \cite{demexico}, enfocado en el diseño, desarrollo preliminar y validación de una solución informática aplicada al ámbito municipal. Por un lado, desde la parte técnica y descriptiva, el estudio detalla la arquitectura del software, el modelado de bases de datos, los componentes de geolocalización (WebGIS) y los requerimientos funcionales de la plataforma web. Por otro lado, en su dimensión propositiva, el proyecto parte de la creación autónoma de un prototipo funcional para orientar la interacción con la Alcaldía Municipal hacia una fase de validación de campo, pruebas de usabilidad y aceptación. Esto garantiza su adaptabilidad operativa de cara a su implementación y defensa, programada para el primer semestre del año 2027. 

\section{Limitaciones}
El desarrollo e implementación de la presente investigación presenta ciertas limitaciones de carácter temporal, institucional y técnico. En primer lugar, en el ámbito temporal, el proyecto se encuentra acotado por el calendario académico de la carrera, estableciendo su defensa para el primer semestre de 2027, lo que delimita los tiempos de desarrollo y evolución del software. En segundo lugar, en el aspecto institucional y de campo, la fase de pruebas y validación del sistema está sujeta a la disponibilidad y a los tiempos de gestión administrativa de la Alcaldía Municipal de Rivas. Finalmente, en el aspecto técnico y operativo, la plataforma está limitada temporalmente por el uso de componentes y servicios cartográficos gratuitos de terceros, los cuales están sujetos a restricciones de cuotas de uso y a la disponibilidad del proveedor externo, lo que podría condicionar el rendimiento del componente WebGIS ante un alto volumen de peticiones simultáneas. 

\section{Viabilidad del Proyecto}
\subsection{Viabilidad Económica}
\subsection{Viabilidad Técnica}
\subsection{Viabilidad Operativa}
\subsection{Viabilidad Legal}
